\section{Data Profiling}

\subsection{Mục tiêu và phạm vi}

\begin{itemize}
    \item \textbf{Dữ liệu đầu vào gồm 4 nguồn:} listings, calendar, reviews, neighbourhoods.geojson.
    \item \textbf{Mục tiêu nghiệp vụ:}
    \begin{itemize}
        \item Task 1: Phân tích tổng quan thị trường và không gian (Market Overview \& Spatial Distribution).
        \item Task 2: Phân tích xu hướng theo thời gian (Temporal Trends).
    \end{itemize}
    \item \textbf{Trọng tâm báo cáo:} Nêu rõ logic đánh giá dữ liệu (profiling) $\rightarrow$ định hướng làm sạch $\rightarrow$ cấu trúc dữ liệu để vẽ biểu đồ.
\end{itemize}

\subsubsection*{1.1 Quy ước chung}

Để bài toán dễ theo dõi, thay vì dùng toàn bộ hàng chục cột dữ liệu gốc, chúng ta chỉ lấy ra các cột thật sự mang ý nghĩa:

\begin{itemize}
    \item \textbf{Với Task 1 (Thị trường \& Không gian):} Cần trả lời câu hỏi ``Bán loại phòng gì, ở đâu, giá bao nhiêu?''. Do đó, \texttt{price} (giá) là thước đo chính, \texttt{room\_type} (loại phòng) để phân nhóm, còn \texttt{neighbourhood} (khu vực) và \texttt{latitude/longitude} (tọa độ) dùng để chấm lên bản đồ.
    \item \textbf{Với Task 2 (Xu hướng thời gian):} Lẽ ra ta sẽ xem giá thay đổi thế nào qua các tháng. Tuy nhiên, vì phát hiện cột giá theo ngày (\texttt{calendar.price}) bị bỏ trống 100\%, nên ta đổi sang quan sát tỷ lệ cư trú. Hai cột \texttt{date} và \texttt{available} (trạng thái phòng) trở thành cốt lõi để vẽ nên biểu đồ mùa vụ.
\end{itemize}

\subsection{Phiên bản dữ liệu (raw $\rightarrow$ selected $\rightarrow$ clean)}

\begin{table}[H]
\centering
\caption{Các phiên bản dữ liệu trong quy trình xử lý}
\begin{tabularx}{\textwidth}{@{} l l X @{}}
\toprule
\textbf{Phiên bản} & \textbf{Ý nghĩa} & \textbf{Dùng cho bước nào} \\
\midrule
raw      & Dữ liệu gốc sau khi nạp file & Kiểm tra hiện trạng và phát hiện lỗi chính \\ \hline
selected & Tập con thuộc tính phục vụ profiling theo task & Đo lường chất lượng trên biến trọng yếu \\ \hline
clean    & Dữ liệu sau pipeline làm sạch và kiểm tra toàn vẹn & Phân tích, tổng hợp KPI, thiết kế biểu đồ \\ \hline
\bottomrule
\end{tabularx}
\end{table}

Quy ước sử dụng trong báo cáo: profiling chủ yếu đọc trên raw/selected để nhận diện vấn đề (đọc ở phần phụ lục); phân tích và dashboard dùng \texttt{*\_clean} để đảm bảo nhất quán và khả năng tái lập.

\subsection{Profiling và chất lượng dữ liệu (tổng quan)}

\subsubsection*{3.1 Quy mô dữ liệu ban đầu}

\begin{table}[H]
\centering
\caption{Quy mô dữ liệu ban đầu}
\begin{tabular}{@{} l r l @{}}
\toprule
\textbf{Dataset} & \textbf{Rows (raw)} & \textbf{Cột profiling chính} \\
\midrule
Listings                    & 36,353         & 13 \\ \hline
Calendar                    & 13,268,858     & 10 \\\hline
Reviews                     & 1,000,870      & 3  \\\hline
GeoJSON (neighbourhoods)    & 233 features   & 2 thuộc tính vùng chính \\
\bottomrule
\end{tabular}
\end{table}

\subsubsection*{3.2 Vấn đề chất lượng quan trọng và tác động phân tích}

Chú thích tra cứu bảng chi tiết (Phụ lục):

\begin{itemize}
    \item \textbf{Kiểm tra tính hợp lệ}: xem Bảng~\ref{tab:A1}.
    \item \textbf{Thống kê mô tả số}: xem Bảng~\ref{tab:A2}.
    \item \textbf{Cấu trúc dữ liệu theo thuộc tính}: xem Bảng~\ref{tab:A3}.
    \item \textbf{Mẫu định dạng}: xem Bảng~\ref{tab:A4}.
    \item \textbf{Hồ sơ ngoại lệ (IQR)}: xem Bảng~\ref{tab:A5}.
    \item \textbf{Bảng giá trị thiếu}: xem Bảng~\ref{tab:A6}.
    \item \textbf{Ví dụ bản ghi lỗi \texttt{minimum\_nights} $>$ \texttt{maximum\_nights}}: xem Bảng~\ref{tab:A7}.
    \item \textbf{\texttt{calendar.price} thiếu 100\% (source-missing):} không thể triển khai xu hướng giá theo thời gian.
    \item \textbf{\texttt{listings.reviews\_per\_month} và \texttt{listings.review\_scores\_rating} thiếu theo ngữ nghĩa (MAR):} gắn với listing chưa có review, không nội suy cưỡng bức.
    \item \textbf{Tồn tại mâu thuẫn logic \texttt{minimum\_nights} $>$ \texttt{maximum\_nights} trong calendar raw:} cần xử lý trước khi tổng hợp theo thời gian.
    \item \textbf{\texttt{maximum\_nights} có giá trị cực đại bất thường (2,147,483,647):} cần kiểm soát khi tóm tắt ràng buộc đặt phòng.
\end{itemize}

Ví dụ lỗi dữ liệu tiêu biểu:

\begin{enumerate}
    \item \textbf{\texttt{calendar.price}}: thiếu toàn bộ $\rightarrow$ loại khỏi mọi phân tích giá theo ngày/tháng.
    \item \textbf{\texttt{minimum\_nights} $>$ \texttt{maximum\_nights}}: biểu hiện mâu thuẫn nghiệp vụ $\rightarrow$ vô hiệu cặp trường nights ở bản ghi lỗi.
\end{enumerate}

\subsubsection*{3.3 Mức sẵn sàng của dữ liệu không gian (GeoJSON)}

Chú thích tra cứu bảng chi tiết (Phụ lục):

\begin{itemize}
    \item \textbf{Tổng quan chỉ số của file GeoJSON}: xem Bảng~\ref{tab:A9}.
    \item \textbf{Kiểm tra tọa độ/bbox}: xem Bảng~\ref{tab:A10}.
    \item \textbf{Coverage join theo neighbourhood}: xem Bảng~\ref{tab:A11}.
\end{itemize}

\begin{itemize}
    \item \texttt{n\_features = 233}, \texttt{n\_neighbourhood = 230}, \texttt{n\_neighbourhood\_group = 5}.
    \item Coverage join dựa theo neighbourhood: match 222, listing-only 0, GeoJSON-only 8.
    \item Ý nghĩa: đủ điều kiện cho point map hoặc choropleth.
\end{itemize}

\subsection{Quy trình làm sạch dữ liệu}

\begin{enumerate}
    \item \textbf{Standardize}: chuẩn hóa kiểu dữ liệu (giá, ngày, nhãn availability) và tạo biến thời gian mới (ngày, tháng).
    \item \textbf{Deduplicate}: loại trùng theo mức độ chi tiết của từng bảng (\texttt{id}, \texttt{listing\_id-date}, review event).
    \item \textbf{Handle missing}: giữ missing có ý nghĩa nghiệp vụ (các trường về reviews), loại hoặc bỏ biến thiếu hoàn toàn (\texttt{calendar.price} với missing 100\%).
    \item \textbf{Remove invalid/inconsistent}: loại bản ghi không hợp lệ, xử lý mâu thuẫn min $>$ max, lọc orphan để đảm bảo 2 bảng hợp lại là 1-n.
    \item \textbf{Outlier policy}: không xóa outlier giá, gắn cờ để phân tích thêm theo bối cảnh thị trường.
\end{enumerate}

Chú thích tra cứu bảng quyết định làm sạch chi tiết: xem Bảng~\ref{tab:A8} (Phụ lục).

\subsection{Tóm tắt dữ liệu sau làm sạch}

\subsubsection*{5.1 Quy mô sau làm sạch}

\begin{table}[H]
\centering
\caption{Quy mô dữ liệu sau làm sạch}
\begin{tabularx}{\textwidth}{@{} l r r X @{}}
\toprule
\textbf{Dataset} & \textbf{Clean rows} & \textbf{Clean cols} & \textbf{Vai trò phân tích} \\
\midrule
listings\_clean  & 21,415    & 14 & Bảng số liệu chính dạng cho thị trường/không gian \\ \hline
calendar\_clean  & 7,816,487 & 11 & Bảng số liệu chính theo thời gian cho tỷ lệ lấp đầy (occupancy) hoặc khả dụng \\ \hline
reviews\_clean   & 787,325   & 3  & Tín hiệu bổ trợ chất lượng/hoạt động \\
\bottomrule
\end{tabularx}
\end{table}

Chú thích tra cứu bảng hậu làm sạch chi tiết (Phụ lục): \textbf{Quy mô giảm theo pipeline} xem Bảng~\ref{tab:A12}; \textbf{Toàn vẹn khóa sau làm sạch} xem Bảng~\ref{tab:A13}.

\subsubsection*{5.2 Thuộc tính trọng yếu được giữ}

\begin{itemize}
    \item \textbf{Không gian:} \texttt{neighbourhood\_group\_cleansed}, \texttt{neighbourhood\_cleansed}, \texttt{latitude}, \texttt{longitude}.
    \item \textbf{Thời gian:} \texttt{date}, \texttt{year}, \texttt{month}.
    \item \textbf{Measure chính:} \texttt{price} (listings), \texttt{is\_available}, \texttt{is\_booked}, \texttt{review\_scores\_rating}, \texttt{reviews\_per\_month}, \texttt{number\_of\_reviews}.
    \item \textbf{Khóa:} \texttt{listings.id}, \texttt{calendar.listing\_id}, \texttt{reviews.listing\_id}.
\end{itemize}

\subsection{Kết luận profiling}

\subsubsection*{6.1 Dữ liệu khả dụng cho phân tích}

\begin{itemize}
    \item \textbf{Khả dụng:} \texttt{listings\_clean}, \texttt{calendar\_clean}, \texttt{reviews\_clean}, \texttt{neighbourhoods.geojson}.
    \item \textbf{Không khả dụng theo mục tiêu ban đầu:} temporal price từ \texttt{calendar.price} do thiếu 100\%.
\end{itemize}

\subsubsection*{6.2 Liên kết chất lượng dữ liệu $\rightarrow$ quyết định phân tích}

\begin{itemize}
    \item Vì \texttt{calendar.price} không dùng được, Task 2 chuyển từ xu hướng giá theo thời gian sang xu hướng lấp đầy/khả dụng (xem Bảng~\ref{tab:A6} và A.8).
    \item Vì missing review fields mang ý nghĩa nghiệp vụ, giữ missing để tránh làm méo diễn giải chất lượng listing (xem Bảng~\ref{tab:A6}, A.8).
    \item Vì có mâu thuẫn và outlier nghiệp vụ, pipeline làm sạch áp dụng nguyên tắc sửa bằng cách đánh dấu thay vì xóa (xem Bảng~\ref{tab:A5}, A.7, A.8).
\end{itemize}

\subsubsection*{6.3 Kết luận phạm vi bài toán}

\begin{itemize}
    \item Giữ Task 1 (market + spatial) trên dữ liệu clean.
    \item Giữ Task 2 theo lấp đầy/khả dụng.
    \item Loại xu hướng giá theo thời gian trên dữ liệu hiện có.
\end{itemize}

Phần tiếp theo chuyển từ bằng chứng chất lượng dữ liệu sang đặc tả mô hình dữ liệu và quy tắc tổng hợp dùng cho dashboard.
